\subsection{2.2.3 Projektive Reduktion und beobachterabhängige Geometrie}

Die bisherigen Untersuchungen führten zu einer grundlegenden Einsicht:

\begin{center}
\emph{Die Geometrie der beobachteten Welt muss nicht identisch mit der Geometrie des OPP sein.}
\end{center}

Vielmehr spricht sowohl die mathematische Struktur der Simulationen als auch die Physik dafür, dass zwischen beiden Ebenen eine Projektionsabbildung existiert.

Formal schreiben wir

\[
\Pi :
\mathcal{O}
\longrightarrow
\mathcal{M},
\]

wobei

\begin{itemize}
\item
\(\mathcal{O}\) den Ontophänomenalen Phasenraum (OPP),

\item
\(\mathcal{M}\) die beobachtbare physikalische Welt,

\item
\(\Pi\) den Projektionsoperator
\end{itemize}

bezeichnet.

Der OPP besitzt keine vorgegebene Raumzeit.

Er besitzt lediglich Relationen.

Erst die Projektion erzeugt eine effektive Geometrie.

\bigskip

Diese Sichtweise besitzt einen bemerkenswerten Vorteil.

Sie erklärt unmittelbar, weshalb frühere numerische Versuche scheiterten.

Denn sämtliche frühen Simulationen gingen implizit davon aus,

\[
\mathcal{M}
=
\mathcal{O},
\]

also

\[
\Pi
=
\mathrm{Id}.
\]

Genau diese Annahme wurde jedoch experimentell widerlegt.

\section*{Der Beobachter als Projektor}

Die Projektionshypothese bedeutet nicht zwangsläufig, dass ein menschlicher Beobachter erforderlich ist.

Der Begriff „Beobachter“ wird hier wesentlich allgemeiner verstanden.

Ein Beobachter ist jedes System, das

\begin{enumerate}
\item Informationen selektiert,
\item Relationen reduziert,
\item eine endliche Auflösung besitzt.
\end{enumerate}

Formal kann eine Projektion daher beschrieben werden als

\[
\Pi
=
\Pi_{\text{Auflösung}}
\circ
\Pi_{\text{Relation}}
\circ
\Pi_{\text{Information}}.
\]

Jede Projektion verliert Information.

Dies ist keine Schwäche,

sondern ihre eigentliche Aufgabe.

\section*{Projektive Kompression}

Eine Projektion reduziert den Informationsgehalt.

Ist

\[
I(\mathcal O)
\]

die Gesamtinformation des OPP,

dann gilt allgemein

\[
I(\Pi(\mathcal O))
<
I(\mathcal O).
\]

Es existiert also keine vollständige Rekonstruktion.

Die beobachtete Physik stellt lediglich einen komprimierten Ausschnitt dar.

Diese Aussage passt erstaunlich gut zu mehreren bekannten Bereichen der modernen Physik:

\begin{itemize}
\item Holographisches Prinzip
\item Renormierungsgruppen
\item Informationsverlust bei Messungen
\item coarse graining
\item effektive Feldtheorien
\end{itemize}

Die PST liefert hierfür einen gemeinsamen Ursprung.

\section*{Warum gerade vier Dimensionen?}

Eine der größten offenen Fragen lautet:

\begin{center}
Warum beobachten wir genau drei Raumdimensionen und eine Zeitdimension?
\end{center}

Die bisherigen numerischen Ergebnisse deuten darauf hin, dass dies keine fundamentale Eigenschaft des OPP ist.

Vielmehr scheint zu gelten

\[
d_{\mathrm{eff}}
=
f(\Pi).
\]

Die Dimension ist also eine Eigenschaft der Projektion.

Nicht des OPP selbst.

Dies ist ein fundamentaler Perspektivwechsel.

Bisherige Theorien fragen:

\begin{quote}
Warum besitzt die Welt vier Dimensionen?
\end{quote}

Die PST fragt stattdessen:

\begin{quote}
Welche Projektion eines hochdimensionalen OPP erzeugt stabil eine vierdimensionale effektive Geometrie?
\end{quote}

\section*{Die Rolle der Diffusionsprojektion}

Besonders bemerkenswert war das Ergebnis der Diffusionsabbildung.

Unter geeigneter Wahl des Diffusionsparameters entstand numerisch

\[
d_{\mathrm{eff}}
\approx
4.
\]

Im Gegensatz zu PCA oder zufälligen Projektionen entstand dieser Wert nicht durch explizite Vorgabe einer Dimension.

Stattdessen entstand er aus einer natürlichen Glättung der relationalen Struktur.

Damit besitzt die Diffusionsprojektion mehrere interessante Eigenschaften.

Sie

\begin{itemize}
\item reduziert lokale Fluktuationen,

\item erhält globale Zusammenhänge,

\item entspricht einem Informationsfluss,

\item besitzt eine natürliche physikalische Interpretation.
\end{itemize}

Gerade dieser letzte Punkt ist bemerkenswert.

Diffusion tritt in nahezu allen Bereichen der Physik auf:

\begin{itemize}
\item Wärmeleitung
\item Brownsche Bewegung
\item Quantenpfadintegrale
\item Renormierungsgruppen
\item Informationsausbreitung
\end{itemize}

Es wäre daher überraschend, wenn gerade diese Projektion keine fundamentale Rolle spielen würde.

\section*{Projektive Quantisierung}

Bereits früh entstand innerhalb des OP die Vermutung,

dass Quantisierung möglicherweise keine fundamentale Eigenschaft der Realität ist.

Die Projektionshypothese liefert hierfür erstmals einen konkreten mathematischen Rahmen.

Sei

\[
x
\in
\mathcal O
\]

ein kontinuierlicher Zustand des OPP.

Die Projektion liefert

\[
\Pi(x)
=
x_i,
\]

wobei

\[
x_i
\]

nun zu einer diskreten Klasse gehört.

Quantisierung entsteht damit nicht durch den OPP,

sondern durch die begrenzte Auflösung der Projektion.

Formal ähnelt dies

\[
\text{Sampling}
\]

oder

\[
\text{Diskretisierung}
\]

kontinuierlicher Daten.

Die Quantenstruktur wäre dann eine Eigenschaft der Beobachtungsebene.

Nicht der ontologischen Ebene.

\section*{Projektive Raumzeit}

Noch weitreichender ist die Konsequenz für die Raumzeit.

Wenn dieselbe Projektion sowohl

\begin{itemize}
\item Quantisierung

als auch

\item Vierdimensionalität
\end{itemize}

erzeugt,

dann besitzen beide möglicherweise dieselbe Ursache.

Dies führt zur Hypothese

\[
\boxed{
\text{Raumzeit und Quantisierung entstehen gemeinsam
durch dieselbe Projektionsabbildung.}
}
\]

Diese Vermutung besitzt mehrere attraktive Eigenschaften.

Sie erklärt gleichzeitig

\begin{enumerate}
\item die endliche Informationsauflösung,
\item diskrete Quantenzahlen,
\item die effektive Vierdimensionalität,
\item die Existenz einer makroskopischen Geometrie.
\end{enumerate}

Damit reduziert sich die Zahl unabhängiger Grundannahmen erheblich.

\section*{Offene Fragen}

Die Projektionshypothese wirft jedoch neue Fragen auf.

Unter anderem:

\begin{enumerate}

\item
Warum bevorzugt die Natur gerade bestimmte Projektionsoperatoren?

\item
Existieren mehrere physikalisch mögliche Projektionen?

\item
Ist die Diffusionsprojektion fundamental oder lediglich eine gute Näherung?

\item
Kann aus geeigneten Projektionsprinzipien automatisch die Lorentz-Invarianz folgen?

\item
Entstehen Eichsymmetrien ebenfalls als Eigenschaften der Projektion?

\end{enumerate}

Diese Fragen bilden den Übergang zu den numerischen Untersuchungen der folgenden Kapitel.

Dort wird untersucht,

welche mathematischen Projektionsoperatoren tatsächlich stabile, physikalisch sinnvolle Strukturen erzeugen und welche lediglich mathematische Artefakte darstellen.

